% !TeX root = ../main.tex
\section{Conclusion}\label{sec:conclusion}
\acresetall

In this paper, we looked into tool support that keeps the data of a \SLR traceable.
We presented \SaiLoR, an open-source desktop tool for screening, data extraction, and consolidation.
A \SaiLoR project consisted of plain JSON files that a schema described and that git versioned.
Reviewers worked independently in their own seats and linked their answers to evidence in the papers.
A consolidator resolved disagreements with the help of agreement statistics, and \SaiLoR merged concurrent changes field by field.
Optional \LLM assistance proposed answers with evidence, but a human reviewed every proposal before \SaiLoR applied it.

We also explained the main architecture decisions.
These were the local files instead of a server, git as collaboration layer, the security model, and the move from Electron to Tauri in version 2.0.0.
To keep the tool itself traceable, we documented 388 requirements.
Of these requirements, 92\% linked to implementing source code and 90\% linked to at least one automated test.
The test suites held 2828 unit tests, 28 integration tests, and 100 Rust tests.
We released \SaiLoR 2.0.0 under the GPL-3.0 license and archived it on Zenodo.

Regarding the tool, we want to port the remaining features of the Electron shell to Tauri and retest the user interface in the web views of all three operating systems.
We also want to verify the local model runtime on \GPUs with CUDA and Vulkan and enable the managed installation once llama.cpp provides the required route.

In future research, we want to evaluate \SaiLoR in a user study with reviewers who conduct a real \SLR.
We want to find out whether the tool reduces the effort of extraction and consolidation and whether the recorded history helps others to trace the review.
We also want to measure how accurate the \LLM proposals are and how they influence the decisions of reviewers.

Lastly, we want to have the requirements reviewed by people other than the developer and check the trace links automatically in the \CI.
